\documentclass[11pt]{article}
\usepackage[margin=1in]{geometry}
\usepackage{amsmath,amssymb}
\usepackage{booktabs}
\usepackage{hyperref}
\usepackage{listings}
\usepackage{xcolor}
\lstset{basicstyle=\ttfamily\small,breaklines=true,frame=single}

\title{Independent Replication --- Poisson Neural Networks\\
       (Jin, Zhang, Kevrekidis, Karniadakis, 2020/2022)}
\author{OpenClaw Subagent \\ LLM judge: Argo \texttt{argo:claude-opus-4.7}}
\date{2026-07-04}

\begin{document}
\maketitle

\begin{abstract}
We independently replicate the core Section~IV-A Lotka--Volterra experiment of
Jin et al., \emph{Learning Poisson systems and trajectories of autonomous systems
via Poisson neural networks} (arXiv:2012.03133, IEEE TNNLS 2022,
DOI \texttt{10.1109/TNNLS.2022.3148734}). Using the authors' reference
implementation unchanged and a plain-MLP baseline of $\sim 16\times$ the
parameter count, we confirm both stable long-time rollouts of PNN and its
superior preservation of the LV invariant, at 15\% of the paper's iteration
budget. Verdict: \textbf{REPLICATED}.
\end{abstract}

\section{Paper summary}
The paper proposes \textbf{Poisson Neural Networks} (PNN) for learning the phase
flow of an arbitrary Poisson system $\dot y = B(y)\nabla H(y)$ --- a superset of
canonical Hamiltonian systems where $B(y) = J^{-1}$. By the Darboux--Lie
theorem, any Poisson system is locally conjugate to canonical Hamiltonian form
by a coordinate transform $\theta$. PNN parameterises
\[
\Phi_{\mathrm{PNN}} \;=\; \theta \circ \varphi_H \circ \theta^{-1},
\qquad \theta \approx \mathrm{INN},\quad \varphi_H \approx \mathrm{SympNet},
\]
where $\theta$ is realised by an Invertible Neural Network (affine-coupling,
closed-form inverse) and $\varphi_H$ is a symplectic-by-construction SympNet
(G-type for LV). The composition is Poisson-preserving by construction.

\section{Claims tested}
\begin{center}
\begin{tabular}{clcc}
\toprule
ID & Claim & Testable via rollouts? & Tested here? \\
\midrule
C1 & Architectural Poisson preservation (Darboux--Lie).      & indirect  & indirect \\
C2 & Stable long-time LV rollouts vs.\ unstructured baseline. & yes       & yes \\
C3 & Small drift of LV invariant $H(u,v)$.                    & yes       & yes \\
C4 & Extended pendulum (odd-dim Poisson).                     & yes       & no \\
C5 & Charged particle, NLS, pixel two-body.                   & yes       & no \\
C6 & Single non-Hamiltonian trajectory (Thm 3).               & yes       & no \\
\bottomrule
\end{tabular}
\end{center}

\section{Method}
\paragraph{Data.} Exactly the Sec.~IV-A LV setup:
$(\dot u,\dot v) = (u(v-2),\, v(1-u))$, three trajectories from
$(1,0.8), (1,1), (1,1.2)$, step $h=0.1$, 100 training points each.
Ground truth generated by the authors' 4th-order Stormer--Verlet integrator in
log-canonical coordinates $(p,q) = (\log u, \log v)$.

\paragraph{Models.}
\begin{itemize}
\item PNN: INN (\texttt{dim=2, split\_dim=1, layers=3, sublayers=2, subwidth=30,
      sigmoid}) $\circ$ G-SympNet (\texttt{dim=2, layers=3, width=30, sigmoid}).
      \textbf{816} parameters.
\item MLP baseline: residual FC net $\mathrm{Linear}(2,64)\!\to\!\tanh\!\to\!
      4\!\times\!(\mathrm{Linear}(64,64)\!\to\!\tanh)\!\to\!\mathrm{Linear}(64,2)$
      with $x_{n+1} = x_n + \mathrm{MLP}(x_n)$. \textbf{12\,802} parameters
      ($\sim 16\times$ PNN).
\end{itemize}

\paragraph{Training.} Adam, $\mathrm{lr}=10^{-3}$, full-batch, 30\,000
iterations for both (15\% of the paper's 200\,000). Wall time on one A100:
PNN 311\,s, MLP 61\,s.

\paragraph{Evaluation.} 1000-step autoregressive rollouts from each of the
three training endpoints; per-step MSE (mean over 3 traj.) and invariant drift
$|H(u_n,v_n)-H(u_0,v_0)|$.

\section{Results}
\subsection{Rollout MSE}
\begin{center}
\begin{tabular}{lrrrrrr}
\toprule
Model & step 100 & step 500 & step 1000 & mean 1--100 & mean 1--500 & mean 1--1000 \\
\midrule
PNN                 & 4.89e-3 & 5.01e-3 & 3.61e-3            & 1.23e-2 & 1.07e-2 & 1.03e-2 \\
MLP ($16\times$ p.) & 4.89e-3 & 2.24e-2 & \textbf{1.63e-1}   & 1.33e-2 & 1.61e-2 & 3.15e-2 \\
ratio MLP/PNN       & 1.0     & 4.5     & \textbf{45}        & ---     & 1.5     & 3.1     \\
\bottomrule
\end{tabular}
\end{center}

\subsection{Invariant drift $|H(u_n,v_n)-H(u_0,v_0)|$}
\begin{center}
\begin{tabular}{lrr}
\toprule
Model & max over 1000 steps & at step 1000 \\
\midrule
PNN                     & 5.81e-3 & 1.41e-3 \\
MLP baseline            & 2.98e-2 & 2.88e-2 \\
Reference (SV, sanity)  & 4.77e-7 & 2.38e-7 \\
\bottomrule
\end{tabular}
\end{center}

PNN's invariant drift is $\sim 5\times$ smaller than the MLP's max and
$\sim 20\times$ smaller at the final step. PNN drift decreases from max to
final; MLP drift saturates near its max.

\section{LLM-judge verdict}
Argo \texttt{argo:claude-opus-4.7} (free tier) returned \texttt{overall :
REPLICATED} with C2 and C3 REPLICATED and C1 OUT-OF-SCOPE (architectural,
not directly testable via rollouts). Full JSON in
\texttt{evidence/judge\_argo.json}.

\section{GENUINE CRITIQUE}
This section records honest limitations that the ``REPLICATED'' verdict does
\emph{not} paper over.

\paragraph{Reduced iteration budget.} Training used 30\,000 iterations
vs.\ the paper's 200\,000 (15\%). Because PNN's rollout MSE at step 1000
(3.61e-3) is already \emph{smaller} than at step 100 (4.89e-3), the model is
clearly in a stable regime; longer training would very likely reduce absolute
error further but is unlikely to overturn the qualitative comparison. Still,
we cannot claim to have matched the paper's exact numeric error level.

\paragraph{Baseline choice.} We compared against a residual MLP with
$\sim 16\times$ more parameters, not against a bare SympNet on multi-trajectory
non-canonical data (the paper's own qualitative counter-example). The MLP
baseline is a fair \emph{unstructured} control but does not directly reproduce
the paper's illustrative failure mode. A follow-up should add the bare-SympNet
baseline the paper implies.

\paragraph{Scope: one experiment out of many.} Only Section IV-A (LV) was
attempted. Claims C4 (extended pendulum, odd-dim Poisson), C5 (charged particle
in EM potential, NLS, pixel two-body), and C6 (single non-Hamiltonian
trajectory, Thm~3) are \emph{unattempted}, not \emph{replicated}. They remain
claims-in-good-standing.

\paragraph{Structural claim (C1) is indirect.} Poisson-structure preservation
by construction is an architectural property of the composition
$\mathrm{INN}\circ\mathrm{SympNet}\circ\mathrm{INN}^{-1}$. Our rollout stability
is \emph{consistent} with it, not a proof of it. The judge correctly marked C1
OUT-OF-SCOPE.

\paragraph{Ground-truth reference is a symplectic integrator, not analytic.}
Errors are measured against SV rollouts of the same ODE, not against a
closed-form solution (LV has no elementary closed form). SV drift is
$\sim 5\times 10^{-7}$, i.e.\ $\sim 4$ orders below PNN error, so this is not
a practical concern.

\paragraph{Single-seed, single-run.} No statistical confidence bands are
reported. The paper is also largely single-seed on Fig.~2, so this matches
paper practice, but a bootstrap over $\ge 5$ seeds would strengthen future
replications.

\paragraph{No test of extrapolation outside the training basin.} Rollouts start
from the training endpoints, which lie \emph{on} the trained level curves.
Extrapolation to fresh initial conditions on nearby but unseen level curves is
untested and would be a stronger stress test of the learned Poisson structure.

\section{Verdict}
\textbf{REPLICATED} on the focal Section IV-A experiment. Author code
(\url{https://github.com/jpzxshi/pnn}) used unchanged; PNN's rollout MSE gap
($\sim 45\times$ at step 1000) and invariant-drift gap ($\sim 5$--$20\times$)
over an $\sim 16\times$-larger unstructured MLP baseline reproduce the paper's
central empirical narrative on Lotka--Volterra. Extensions to
C4--C6 remain unattempted and unclaimed.

\end{document}
